\section{Conclusion}
\label{sec:conclusion}

We have characterised induction circuit stability under fine-tuning in a
two-layer attention-only transformer, measuring per-head induction scores
and activation-patching attribution at every checkpoint across both a Python
code condition and a TinyStories prose control.

\textbf{Summary.} The pretrained \texttt{attn-only-2l} model contains a
single-head induction circuit (L1H6: prefix-matching score $0.408 \pm 0.103$,
attribution $0.952$, $12.4\times$ the next head). Under both fine-tuning
conditions, this circuit \emph{strengthens} rather than degrades, reaching
$0.638$ (code) and $0.616$ (prose) after 244 training steps. Trajectories
are near-identical across domains; both exceed the $0.5$ IS threshold by
step 100 and continue rising. Adversarial probes confirm the post-fine-tuning
head is a stronger, selective induction mechanism: structure-preserving
perturbations retain $0.944$ of the clean score; structure-removing probes
drop to $0.132$. The degenerate \texttt{all\_same\_token} outlier decreases
from $4.12$ to $2.41$ after fine-tuning.

The interactive visualisation dashboard and full codebase are released
alongside this paper, enabling practitioners to apply circuit-level
monitoring to their own fine-tuning pipelines.

\textbf{Limitations of current results.} Per-head trajectories (L1H6
specifically) are from seed 42 only. Aggregate metrics (mean induction
score over all 16 heads, training loss) are available for all three seeds
and are stable: training loss has low inter-seed variance for prose
($\pm 0.08$) and higher variance for code ($\pm 0.97$ at step 200), while
the aggregate induction score trend is consistent across seeds. Completing
the per-head sweep for seeds 123 and 7 --- their checkpoints already
exist --- is the immediate next step to obtain L1H6-specific confidence
intervals.

\textbf{Future work.}
(1)~Complete the 3-seed sweep and add $\pm 1$ SD shading to Figures~5
and~6 to confirm the reinforcement direction with statistical precision.
(2)~Extend the token budget to test whether the score plateau between
steps 100 and 200 continues, reverses, or accelerates at larger scale.
(3)~Investigate whether the observed IS increase under code fine-tuning
increases susceptibility to induction-based prompt injection, and whether
it can be controlled via circuit-targeted regularisation.
(4)~Extend to GPT-2 small (with MLP layers) using automated circuit
discovery~\citep{conmy2023automated} to test whether the same reinforcement
pattern holds in a model with non-attention computation.
